\documentclass[11pt, a4paper]{article}\usepackage[]{graphicx}\usepackage[]{xcolor}
% maxwidth is the original width if it is less than linewidth
% otherwise use linewidth (to make sure the graphics do not exceed the margin)
\makeatletter
\def\maxwidth{ %
  \ifdim\Gin@nat@width>\linewidth
    \linewidth
  \else
    \Gin@nat@width
  \fi
}
\makeatother

\definecolor{fgcolor}{rgb}{0.345, 0.345, 0.345}
\newcommand{\hlnum}[1]{\textcolor[rgb]{0.686,0.059,0.569}{#1}}%
\newcommand{\hlsng}[1]{\textcolor[rgb]{0.192,0.494,0.8}{#1}}%
\newcommand{\hlcom}[1]{\textcolor[rgb]{0.678,0.584,0.686}{\textit{#1}}}%
\newcommand{\hlopt}[1]{\textcolor[rgb]{0,0,0}{#1}}%
\newcommand{\hldef}[1]{\textcolor[rgb]{0.345,0.345,0.345}{#1}}%
\newcommand{\hlkwa}[1]{\textcolor[rgb]{0.161,0.373,0.58}{\textbf{#1}}}%
\newcommand{\hlkwb}[1]{\textcolor[rgb]{0.69,0.353,0.396}{#1}}%
\newcommand{\hlkwc}[1]{\textcolor[rgb]{0.333,0.667,0.333}{#1}}%
\newcommand{\hlkwd}[1]{\textcolor[rgb]{0.737,0.353,0.396}{\textbf{#1}}}%
\let\hlipl\hlkwb

\usepackage{framed}
\makeatletter
\newenvironment{kframe}{%
 \def\at@end@of@kframe{}%
 \ifinner\ifhmode%
  \def\at@end@of@kframe{\end{minipage}}%
  \begin{minipage}{\columnwidth}%
 \fi\fi%
 \def\FrameCommand##1{\hskip\@totalleftmargin \hskip-\fboxsep
 \colorbox{shadecolor}{##1}\hskip-\fboxsep
     % There is no \\@totalrightmargin, so:
     \hskip-\linewidth \hskip-\@totalleftmargin \hskip\columnwidth}%
 \MakeFramed {\advance\hsize-\width
   \@totalleftmargin\z@ \linewidth\hsize
   \@setminipage}}%
 {\par\unskip\endMakeFramed%
 \at@end@of@kframe}
\makeatother

\definecolor{shadecolor}{rgb}{.97, .97, .97}
\definecolor{messagecolor}{rgb}{0, 0, 0}
\definecolor{warningcolor}{rgb}{1, 0, 1}
\definecolor{errorcolor}{rgb}{1, 0, 0}
\newenvironment{knitrout}{}{} % an empty environment to be redefined in TeX

\usepackage{alltt}

\usepackage[top = 0.6 in, bottom = 0.6 in, left = 0.8 in, right = 0.8 in]{geometry}
\usepackage{amsmath, amssymb, amsfonts}

\allowdisplaybreaks[4]

\usepackage{enumerate}
\usepackage{array}
\usepackage{multirow}
\usepackage{dingbat}
\usepackage{fontawesome5}
\usepackage{tasks}
\usepackage{bbding}
\usepackage{undertilde}
\usepackage{twemojis}
% how to use bull's eye ----- \scalebox{2.0}{\twemoji{bullseye}}
\usepackage{fontspec}
\usepackage{fancyhdr}
\usepackage{lastpage}

\pagestyle{fancy}
\fancyhf{}
\fancyfoot[C]{Page \thepage\ of \pageref{LastPage}}
\renewcommand{\headrulewidth}{0pt}

\title{MSMS 408 : Practical 14}
\author{{\fontsize{25}{20} \benyorithafont Ananda Biswas} \\[1em] Exam Roll No. : 24419STC053}
\date{\today}

\newfontface\benyorithafont{Benyoritha-G334D.otf}

\newfontface\myfont{Myfont1-Regular.ttf}[LetterSpace=0.05em]

\newfontface\cbfont{CaveatBrush-Regular.ttf}
% how to use --- \myfont --write text here--

\newfontface\gloriahallelujahfont{GLORIAHALLELUJAH-REGULAR.TTF}
\IfFileExists{upquote.sty}{\usepackage{upquote}}{}
\begin{document}

\maketitle


\section*{\faArrowAltCircleRight[regular] \textcolor{blue}{Question}}

\settasks{label = (\roman*),
          label-width = 4ex,
          item-indent = 8em}

\hspace{1cm} Given the functions
\begin{tasks}(3)
\task $\displaystyle{f(x) = (x - 1)^2}$,
\task $\displaystyle{f(x) = e^{-x}}$,
\task $\displaystyle{f(x) = x^2}$
\end{tasks}

each defined on $[0, 2]$


\begin{enumerate}[(i)]
\item examine their monotonicity using first derivatives.

\item For the monotonic functions, estimate $E[f(X)]$ using the Monte Carlo method.

\item For the monotonic functions, estimate $E[f(X)]$ using the Antithetic variable technique.

\item Compute and compare the variances of both the estimates for sample sizes $30, 80, 100, 150$.
\end{enumerate}

\section*{\faArrowAltCircleRight[regular] \textcolor{blue}{Theory}}

\begin{enumerate}

\item[\scalebox{2.0}{\twemoji{keycap: 1}}] 

\begin{enumerate}
\item $f(x) = (x - 1)^2 \Rightarrow f'(x) = 2(x - 1)$. $f(x)$ is decreasing when $x \in [0, 1)$ and $f(x)$ is increasing when $x \in (1, 2]$. So $f(x)$ is not monotonic in $[0, 2]$.

\vspace{1em}

\item $f(x) = e^{-x} \Rightarrow f'(x) = - e^{-x} < 0 \; \forall x \in [0, 2]$. So $f(x)$ is monotonically decreasing in $[0, 2]$.

\vspace{1em}

\item $f(x) = x^2 \Rightarrow f'(x) = 2x \geq 0 \; \forall x \in [0, 2]$. So $f(x)$ is monotonically increasing in $[0, 2]$.
\end{enumerate}

\vspace{1em}

$\bullet$ If $x_1, x_2, \ldots, x_n$ is a sample of size $n$ from $X$, Monte Carlo estimate of $\theta = E[f(X)]$ is $$\widehat{\theta} = \dfrac{1}{n} \sum \limits_{i = 1}^{n} f(x_i).$$

$\bullet$ Estimate of $\theta$ by using antithetic variable will be

$$\widehat{\theta} = \dfrac{1}{n} \sum \limits_{i = 1}^{n/2} [f(X_i) + f(X_i')]$$

where $X_i$ and $X_i'$ are negatively correlated random variables. \\[0.5em]

Here $h(u) = 2u$.

\end{enumerate}

\newpage

\section*{\faArrowAltCircleRight[regular] \textcolor{blue}{R Program}}

\begin{enumerate}

\item[\scalebox{2.0}{\twemoji{keycap: 2}}]





\begin{knitrout}
\definecolor{shadecolor}{rgb}{0.969, 0.969, 0.969}\color{fgcolor}\begin{kframe}
\begin{alltt}
\hldef{h} \hlkwb{<-} \hlkwa{function}\hldef{(}\hlkwc{u}\hldef{)} \hlnum{2} \hlopt{*} \hldef{u}
\end{alltt}
\end{kframe}
\end{knitrout}

\begin{knitrout}
\definecolor{shadecolor}{rgb}{0.969, 0.969, 0.969}\color{fgcolor}\begin{kframe}
\begin{alltt}
\hldef{sample.sizes} \hlkwb{<-} \hlkwd{c}\hldef{(}\hlnum{30}\hldef{,} \hlnum{80}\hldef{,} \hlnum{100}\hldef{,} \hlnum{150}\hldef{)}
\end{alltt}
\end{kframe}
\end{knitrout}

\begin{knitrout}
\definecolor{shadecolor}{rgb}{0.969, 0.969, 0.969}\color{fgcolor}\begin{kframe}
\begin{alltt}
\hldef{f1} \hlkwb{<-} \hlkwa{function}\hldef{(}\hlkwc{x}\hldef{)} \hlkwd{exp}\hldef{(}\hlopt{-}\hldef{x)}
\end{alltt}
\end{kframe}
\end{knitrout}

\begin{knitrout}
\definecolor{shadecolor}{rgb}{0.969, 0.969, 0.969}\color{fgcolor}\begin{kframe}
\begin{alltt}
\hldef{f1.est.mc} \hlkwb{<-} \hlkwd{c}\hldef{()}

\hlkwa{for} \hldef{(i} \hlkwa{in} \hlnum{1}\hlopt{:}\hlkwd{length}\hldef{(sample.sizes)) \{}
  \hldef{size} \hlkwb{<-} \hldef{sample.sizes[i]}

  \hldef{u} \hlkwb{<-} \hlkwd{runif}\hldef{(size)}

  \hldef{x} \hlkwb{<-} \hlkwd{h}\hldef{(u)}

  \hldef{f1.est.mc[i]} \hlkwb{<-} \hlkwd{mean}\hldef{(}\hlkwd{f1}\hldef{(x))}
\hldef{\}}
\end{alltt}
\end{kframe}
\end{knitrout}

\begin{knitrout}
\definecolor{shadecolor}{rgb}{0.969, 0.969, 0.969}\color{fgcolor}\begin{kframe}
\begin{alltt}
\hldef{df1} \hlkwb{<-} \hlkwd{data.frame}\hldef{(}\hlkwc{sample.size} \hldef{= sample.sizes,}
                  \hlkwc{estimate} \hldef{= f1.est.mc)}
\end{alltt}
\end{kframe}
\end{knitrout}




% Table created by stargazer v.5.2.3 by Marek Hlavac, Social Policy Institute. E-mail: marek.hlavac at gmail.com
% Date and time: Tue, Apr 07, 2026 - 22:28:25
\begin{table}[!htbp] \centering 
  \caption{Monte Carlo estimate of E[f1(X)]} 
  \label{Table 1} 
\begin{tabular}{@{\extracolsep{5pt}} cc} 
\\[-1.8ex]\hline 
\hline \\[-1.8ex] 
sample.size & estimate \\ 
\hline \\[-1.8ex] 
$30$ & $0.343139$ \\ 
$80$ & $0.453867$ \\ 
$100$ & $0.427947$ \\ 
$150$ & $0.438943$ \\ 
\hline \\[-1.8ex] 
\end{tabular} 
\end{table} 


The true value is $E[e^{-X}] = \displaystyle{\int \limits_{0}^{2} e^{-x} \cdot \dfrac{1}{2} dx} = \dfrac{1}{2} (1 - e^{-2}) = 0.4323324$.

\begin{knitrout}
\definecolor{shadecolor}{rgb}{0.969, 0.969, 0.969}\color{fgcolor}\begin{kframe}
\begin{alltt}
\hldef{f2} \hlkwb{<-} \hlkwa{function}\hldef{(}\hlkwc{x}\hldef{) x}\hlopt{^}\hlnum{2}
\end{alltt}
\end{kframe}
\end{knitrout}

\begin{knitrout}
\definecolor{shadecolor}{rgb}{0.969, 0.969, 0.969}\color{fgcolor}\begin{kframe}
\begin{alltt}
\hldef{f2.est.mc} \hlkwb{<-} \hlkwd{c}\hldef{()}

\hlkwa{for} \hldef{(i} \hlkwa{in} \hlnum{1}\hlopt{:}\hlkwd{length}\hldef{(sample.sizes)) \{}
  \hldef{size} \hlkwb{<-} \hldef{sample.sizes[i]}

  \hldef{u} \hlkwb{<-} \hlkwd{runif}\hldef{(size)}

  \hldef{x} \hlkwb{<-} \hlkwd{h}\hldef{(u)}

  \hldef{f2.est.mc[i]} \hlkwb{<-} \hlkwd{mean}\hldef{(}\hlkwd{f2}\hldef{(x))}
\hldef{\}}
\end{alltt}
\end{kframe}
\end{knitrout}

\begin{knitrout}
\definecolor{shadecolor}{rgb}{0.969, 0.969, 0.969}\color{fgcolor}\begin{kframe}
\begin{alltt}
\hldef{df2} \hlkwb{<-} \hlkwd{data.frame}\hldef{(}\hlkwc{sample.size} \hldef{= sample.sizes,}
                  \hlkwc{estimate} \hldef{= f2.est.mc)}
\end{alltt}
\end{kframe}
\end{knitrout}


% Table created by stargazer v.5.2.3 by Marek Hlavac, Social Policy Institute. E-mail: marek.hlavac at gmail.com
% Date and time: Tue, Apr 07, 2026 - 22:28:26
\begin{table}[!htbp] \centering 
  \caption{Monte Carlo estimate of E[f2(X)]} 
  \label{Table 2} 
\begin{tabular}{@{\extracolsep{5pt}} cc} 
\\[-1.8ex]\hline 
\hline \\[-1.8ex] 
sample.size & estimate \\ 
\hline \\[-1.8ex] 
$30$ & $1.325422$ \\ 
$80$ & $1.369527$ \\ 
$100$ & $1.289178$ \\ 
$150$ & $1.207762$ \\ 
\hline \\[-1.8ex] 
\end{tabular} 
\end{table} 


The true value is $E[X^{2}] = \displaystyle{\int \limits_{0}^{2} x^{2} \cdot \dfrac{1}{2} dx} = \dfrac{4}{3} = 1.3333333$.

\item[\scalebox{2.0}{\twemoji{keycap: 3}}]

\begin{knitrout}
\definecolor{shadecolor}{rgb}{0.969, 0.969, 0.969}\color{fgcolor}\begin{kframe}
\begin{alltt}
\hldef{f1.est.antithetic} \hlkwb{<-} \hlkwd{c}\hldef{()}

\hlkwa{for} \hldef{(i} \hlkwa{in} \hlnum{1}\hlopt{:}\hlkwd{length}\hldef{(sample.sizes)) \{}
  \hldef{size} \hlkwb{<-} \hldef{sample.sizes[i]}

  \hldef{u} \hlkwb{<-} \hlkwd{runif}\hldef{(size}\hlopt{/}\hlnum{2}\hldef{)}

  \hldef{x} \hlkwb{<-} \hlkwd{h}\hldef{(u)}
  \hldef{x_dash} \hlkwb{<-} \hlkwd{h}\hldef{(}\hlnum{1} \hlopt{-} \hldef{u)}

  \hldef{f1.est.antithetic[i]} \hlkwb{<-} \hlkwd{mean}\hldef{(}\hlkwd{c}\hldef{(}\hlkwd{f1}\hldef{(x),} \hlkwd{f1}\hldef{(x_dash)))}
\hldef{\}}
\end{alltt}
\end{kframe}
\end{knitrout}

\begin{knitrout}
\definecolor{shadecolor}{rgb}{0.969, 0.969, 0.969}\color{fgcolor}\begin{kframe}
\begin{alltt}
\hldef{df3} \hlkwb{<-} \hlkwd{data.frame}\hldef{(}\hlkwc{sample.size} \hldef{= sample.sizes,}
                  \hlkwc{estimate} \hldef{= f1.est.antithetic)}
\end{alltt}
\end{kframe}
\end{knitrout}


% Table created by stargazer v.5.2.3 by Marek Hlavac, Social Policy Institute. E-mail: marek.hlavac at gmail.com
% Date and time: Tue, Apr 07, 2026 - 22:28:26
\begin{table}[!htbp] \centering 
  \caption{Estimate of E[f1(X)] from Antithetic Variable Technique} 
  \label{Table 3} 
\begin{tabular}{@{\extracolsep{5pt}} cc} 
\\[-1.8ex]\hline 
\hline \\[-1.8ex] 
sample.size & estimate \\ 
\hline \\[-1.8ex] 
$30$ & $0.427867$ \\ 
$80$ & $0.425985$ \\ 
$100$ & $0.452750$ \\ 
$150$ & $0.438812$ \\ 
\hline \\[-1.8ex] 
\end{tabular} 
\end{table} 


\begin{knitrout}
\definecolor{shadecolor}{rgb}{0.969, 0.969, 0.969}\color{fgcolor}\begin{kframe}
\begin{alltt}
\hldef{f2} \hlkwb{<-} \hlkwa{function}\hldef{(}\hlkwc{x}\hldef{) x}\hlopt{^}\hlnum{2}
\end{alltt}
\end{kframe}
\end{knitrout}

\begin{knitrout}
\definecolor{shadecolor}{rgb}{0.969, 0.969, 0.969}\color{fgcolor}\begin{kframe}
\begin{alltt}
\hldef{f2.est.antithetic} \hlkwb{<-} \hlkwd{c}\hldef{()}

\hlkwa{for} \hldef{(i} \hlkwa{in} \hlnum{1}\hlopt{:}\hlkwd{length}\hldef{(sample.sizes)) \{}
  \hldef{size} \hlkwb{<-} \hldef{sample.sizes[i]}

  \hldef{u} \hlkwb{<-} \hlkwd{runif}\hldef{(size}\hlopt{/}\hlnum{2}\hldef{)}

  \hldef{x} \hlkwb{<-} \hlkwd{h}\hldef{(u)}
  \hldef{x_dash} \hlkwb{<-} \hlkwd{h}\hldef{(}\hlnum{1} \hlopt{-} \hldef{u)}

  \hldef{f2.est.antithetic[i]} \hlkwb{<-} \hlkwd{mean}\hldef{(}\hlkwd{c}\hldef{(}\hlkwd{f2}\hldef{(x),} \hlkwd{f2}\hldef{(x_dash)))}
\hldef{\}}
\end{alltt}
\end{kframe}
\end{knitrout}

\begin{knitrout}
\definecolor{shadecolor}{rgb}{0.969, 0.969, 0.969}\color{fgcolor}\begin{kframe}
\begin{alltt}
\hldef{df4} \hlkwb{<-} \hlkwd{data.frame}\hldef{(}\hlkwc{sample.size} \hldef{= sample.sizes,}
                  \hlkwc{estimate} \hldef{= f2.est.antithetic)}
\end{alltt}
\end{kframe}
\end{knitrout}



% Table created by stargazer v.5.2.3 by Marek Hlavac, Social Policy Institute. E-mail: marek.hlavac at gmail.com
% Date and time: Tue, Apr 07, 2026 - 22:28:26
\begin{table}[!htbp] \centering 
  \caption{Estimate of E[f2(X)] from Antithetic Variable Technique} 
  \label{Table 4} 
\begin{tabular}{@{\extracolsep{5pt}} cc} 
\\[-1.8ex]\hline 
\hline \\[-1.8ex] 
sample.size & estimate \\ 
\hline \\[-1.8ex] 
$30$ & $1.287661$ \\ 
$80$ & $1.394847$ \\ 
$100$ & $1.352673$ \\ 
$150$ & $1.358554$ \\ 
\hline \\[-1.8ex] 
\end{tabular} 
\end{table} 


\item[\scalebox{2.0}{\twemoji{keycap: 4}}]

\begin{knitrout}
\definecolor{shadecolor}{rgb}{0.969, 0.969, 0.969}\color{fgcolor}\begin{kframe}
\begin{alltt}
\hldef{B} \hlkwb{<-} \hlnum{1000}
\end{alltt}
\end{kframe}
\end{knitrout}

\begin{knitrout}
\definecolor{shadecolor}{rgb}{0.969, 0.969, 0.969}\color{fgcolor}\begin{kframe}
\begin{alltt}
\hldef{var.f1.mc} \hlkwb{<-} \hlkwd{c}\hldef{()}

\hlkwa{for} \hldef{(j} \hlkwa{in} \hlnum{1}\hlopt{:}\hlkwd{length}\hldef{(sample.sizes)) \{}

  \hldef{size} \hlkwb{<-} \hldef{sample.sizes[j]}

  \hldef{f1.est.mc} \hlkwb{<-} \hlkwd{c}\hldef{()}

  \hlkwa{for} \hldef{(i} \hlkwa{in} \hlnum{1}\hlopt{:}\hldef{B) \{}
    \hldef{u} \hlkwb{<-} \hlkwd{runif}\hldef{(size)}

    \hldef{x} \hlkwb{<-} \hlkwd{h}\hldef{(u)}

    \hldef{f1.est.mc[i]} \hlkwb{<-} \hlkwd{mean}\hldef{(}\hlkwd{f1}\hldef{(x))}
  \hldef{\}}

  \hldef{var.f1.mc[j]} \hlkwb{<-} \hlkwd{var}\hldef{(f1.est.mc)}
\hldef{\}}
\end{alltt}
\end{kframe}
\end{knitrout}

\begin{knitrout}
\definecolor{shadecolor}{rgb}{0.969, 0.969, 0.969}\color{fgcolor}\begin{kframe}
\begin{alltt}
\hldef{df5} \hlkwb{<-} \hlkwd{data.frame}\hldef{(}\hlkwc{sample.size} \hldef{= sample.sizes,}
                  \hlkwc{variances} \hldef{= var.f1.mc)}
\end{alltt}
\end{kframe}
\end{knitrout}


% Table created by stargazer v.5.2.3 by Marek Hlavac, Social Policy Institute. E-mail: marek.hlavac at gmail.com
% Date and time: Tue, Apr 07, 2026 - 22:28:26
\begin{table}[!htbp] \centering 
  \caption{Variance of estimate of E[f1(X)] from Monte Carlo Method} 
  \label{Table 5} 
\begin{tabular}{@{\extracolsep{5pt}} cc} 
\\[-1.8ex]\hline 
\hline \\[-1.8ex] 
sample.size & variances \\ 
\hline \\[-1.8ex] 
$30$ & $0.001918$ \\ 
$80$ & $0.000675$ \\ 
$100$ & $0.000568$ \\ 
$150$ & $0.000365$ \\ 
\hline \\[-1.8ex] 
\end{tabular} 
\end{table} 


\newpage

\begin{knitrout}
\definecolor{shadecolor}{rgb}{0.969, 0.969, 0.969}\color{fgcolor}\begin{kframe}
\begin{alltt}
\hldef{var.f1.antithetic} \hlkwb{<-} \hlkwd{c}\hldef{()}

\hlkwa{for} \hldef{(j} \hlkwa{in} \hlnum{1}\hlopt{:}\hlkwd{length}\hldef{(sample.sizes)) \{}

  \hldef{size} \hlkwb{<-} \hldef{sample.sizes[j]}

  \hldef{f1.est.antithetic} \hlkwb{<-} \hlkwd{c}\hldef{()}

  \hlkwa{for} \hldef{(i} \hlkwa{in} \hlnum{1}\hlopt{:}\hldef{B) \{}
    \hldef{u} \hlkwb{<-} \hlkwd{runif}\hldef{(size}\hlopt{/}\hlnum{2}\hldef{)}

    \hldef{x} \hlkwb{<-} \hlkwd{h}\hldef{(u)}
    \hldef{x_dash} \hlkwb{<-} \hlkwd{h}\hldef{(}\hlnum{1} \hlopt{-} \hldef{u)}

    \hldef{f1.est.antithetic[i]} \hlkwb{<-} \hlkwd{mean}\hldef{(}\hlkwd{c}\hldef{(}\hlkwd{f1}\hldef{(x),} \hlkwd{f1}\hldef{(x_dash)))}
  \hldef{\}}

  \hldef{var.f1.antithetic[j]} \hlkwb{<-} \hlkwd{var}\hldef{(f1.est.antithetic)}
\hldef{\}}
\end{alltt}
\end{kframe}
\end{knitrout}

\begin{knitrout}
\definecolor{shadecolor}{rgb}{0.969, 0.969, 0.969}\color{fgcolor}\begin{kframe}
\begin{alltt}
\hldef{df6} \hlkwb{<-} \hlkwd{data.frame}\hldef{(}\hlkwc{sample.size} \hldef{= sample.sizes,}
                  \hlkwc{variances} \hldef{= var.f1.antithetic)}
\end{alltt}
\end{kframe}
\end{knitrout}


% Table created by stargazer v.5.2.3 by Marek Hlavac, Social Policy Institute. E-mail: marek.hlavac at gmail.com
% Date and time: Tue, Apr 07, 2026 - 22:28:26
\begin{table}[!htbp] \centering 
  \caption{Variance of estimate of E[f1(X)] from Antithetic Variable Technique} 
  \label{Table 6} 
\begin{tabular}{@{\extracolsep{5pt}} cc} 
\\[-1.8ex]\hline 
\hline \\[-1.8ex] 
sample.size & variances \\ 
\hline \\[-1.8ex] 
$30$ & $0.000249$ \\ 
$80$ & $0.000085$ \\ 
$100$ & $0.000074$ \\ 
$150$ & $0.000050$ \\ 
\hline \\[-1.8ex] 
\end{tabular} 
\end{table} 


\begin{knitrout}
\definecolor{shadecolor}{rgb}{0.969, 0.969, 0.969}\color{fgcolor}\begin{kframe}
\begin{alltt}
\hldef{relative.gain.f1} \hlkwb{<-} \hldef{((var.f1.mc} \hlopt{-} \hldef{var.f1.antithetic)} \hlopt{/} \hldef{var.f1.mc)} \hlopt{*} \hlnum{100}
\end{alltt}
\end{kframe}
\end{knitrout}

\begin{knitrout}
\definecolor{shadecolor}{rgb}{0.969, 0.969, 0.969}\color{fgcolor}\begin{kframe}
\begin{alltt}
\hldef{df7} \hlkwb{<-} \hlkwd{data.frame}\hldef{(}\hlkwc{sample.size} \hldef{= sample.sizes,} \hlkwc{gain} \hldef{= relative.gain.f1)}
\end{alltt}
\end{kframe}
\end{knitrout}


% Table created by stargazer v.5.2.3 by Marek Hlavac, Social Policy Institute. E-mail: marek.hlavac at gmail.com
% Date and time: Tue, Apr 07, 2026 - 22:28:26
\begin{table}[!htbp] \centering 
  \caption{Percentage Variance Reduction in estimation of E[f1(X)]} 
  \label{Table 7} 
\begin{tabular}{@{\extracolsep{5pt}} cc} 
\\[-1.8ex]\hline 
\hline \\[-1.8ex] 
sample.size & gain \\ 
\hline \\[-1.8ex] 
$30$ & $87.032$ \\ 
$80$ & $87.345$ \\ 
$100$ & $86.947$ \\ 
$150$ & $86.302$ \\ 
\hline \\[-1.8ex] 
\end{tabular} 
\end{table} 


\newpage

\smallpencil \hspace{0.1cm} {\setlength{\spaceskip}{1em plus 0.5em minus 0.5em} \fontsize{17}{20}\myfont We achieve significant reduction in variance of the estimates by using antithetic variable technique. \par}

\begin{knitrout}
\definecolor{shadecolor}{rgb}{0.969, 0.969, 0.969}\color{fgcolor}\begin{kframe}
\begin{alltt}
\hldef{var.f2.mc} \hlkwb{<-} \hlkwd{c}\hldef{()}

\hlkwa{for} \hldef{(j} \hlkwa{in} \hlnum{1}\hlopt{:}\hlkwd{length}\hldef{(sample.sizes)) \{}

  \hldef{size} \hlkwb{<-} \hldef{sample.sizes[j]}

  \hldef{f2.est.mc} \hlkwb{<-} \hlkwd{c}\hldef{()}

  \hlkwa{for} \hldef{(i} \hlkwa{in} \hlnum{1}\hlopt{:}\hldef{B) \{}
    \hldef{u} \hlkwb{<-} \hlkwd{runif}\hldef{(size)}

    \hldef{x} \hlkwb{<-} \hlkwd{h}\hldef{(u)}

    \hldef{f2.est.mc[i]} \hlkwb{<-} \hlkwd{mean}\hldef{(}\hlkwd{f2}\hldef{(x))}
  \hldef{\}}

  \hldef{var.f2.mc[j]} \hlkwb{<-} \hlkwd{var}\hldef{(f2.est.mc)}
\hldef{\}}
\end{alltt}
\end{kframe}
\end{knitrout}

\begin{knitrout}
\definecolor{shadecolor}{rgb}{0.969, 0.969, 0.969}\color{fgcolor}\begin{kframe}
\begin{alltt}
\hldef{df8} \hlkwb{<-} \hlkwd{data.frame}\hldef{(}\hlkwc{sample.size} \hldef{= sample.sizes,}
                  \hlkwc{variances} \hldef{= var.f2.mc)}
\end{alltt}
\end{kframe}
\end{knitrout}


% Table created by stargazer v.5.2.3 by Marek Hlavac, Social Policy Institute. E-mail: marek.hlavac at gmail.com
% Date and time: Tue, Apr 07, 2026 - 22:28:27
\begin{table}[!htbp] \centering 
  \caption{Variance of estimate of E[f2(X)] from Antithetic Variable Method} 
  \label{Table 8} 
\begin{tabular}{@{\extracolsep{5pt}} cc} 
\\[-1.8ex]\hline 
\hline \\[-1.8ex] 
sample.size & variances \\ 
\hline \\[-1.8ex] 
$30$ & $0.044867$ \\ 
$80$ & $0.018023$ \\ 
$100$ & $0.013553$ \\ 
$150$ & $0.009031$ \\ 
\hline \\[-1.8ex] 
\end{tabular} 
\end{table} 


\newpage

\begin{knitrout}
\definecolor{shadecolor}{rgb}{0.969, 0.969, 0.969}\color{fgcolor}\begin{kframe}
\begin{alltt}
\hldef{var.f2.antithetic} \hlkwb{<-} \hlkwd{c}\hldef{()}

\hlkwa{for} \hldef{(j} \hlkwa{in} \hlnum{1}\hlopt{:}\hlkwd{length}\hldef{(sample.sizes)) \{}

  \hldef{size} \hlkwb{<-} \hldef{sample.sizes[j]}

  \hldef{f2.est.antithetic} \hlkwb{<-} \hlkwd{c}\hldef{()}

  \hlkwa{for} \hldef{(i} \hlkwa{in} \hlnum{1}\hlopt{:}\hldef{B) \{}
    \hldef{u} \hlkwb{<-} \hlkwd{runif}\hldef{(size}\hlopt{/}\hlnum{2}\hldef{)}

    \hldef{x} \hlkwb{<-} \hlkwd{h}\hldef{(u)}
    \hldef{x_dash} \hlkwb{<-} \hlkwd{h}\hldef{(}\hlnum{1} \hlopt{-} \hldef{u)}

    \hldef{f2.est.antithetic[i]} \hlkwb{<-} \hlkwd{mean}\hldef{(}\hlkwd{c}\hldef{(}\hlkwd{f2}\hldef{(x),} \hlkwd{f2}\hldef{(x_dash)))}
  \hldef{\}}

  \hldef{var.f2.antithetic[j]} \hlkwb{<-} \hlkwd{var}\hldef{(f2.est.antithetic)}
\hldef{\}}
\end{alltt}
\end{kframe}
\end{knitrout}

\begin{knitrout}
\definecolor{shadecolor}{rgb}{0.969, 0.969, 0.969}\color{fgcolor}\begin{kframe}
\begin{alltt}
\hldef{df9} \hlkwb{<-} \hlkwd{data.frame}\hldef{(}\hlkwc{sample.size} \hldef{= sample.sizes,}
                  \hlkwc{variances} \hldef{= var.f2.antithetic)}
\end{alltt}
\end{kframe}
\end{knitrout}


% Table created by stargazer v.5.2.3 by Marek Hlavac, Social Policy Institute. E-mail: marek.hlavac at gmail.com
% Date and time: Tue, Apr 07, 2026 - 22:28:27
\begin{table}[!htbp] \centering 
  \caption{Variance of estimate of E[f2(X)] from Antithetic Variable Technique} 
  \label{Table 9} 
\begin{tabular}{@{\extracolsep{5pt}} cc} 
\\[-1.8ex]\hline 
\hline \\[-1.8ex] 
sample.size & variances \\ 
\hline \\[-1.8ex] 
$30$ & $0.005751$ \\ 
$80$ & $0.002204$ \\ 
$100$ & $0.001778$ \\ 
$150$ & $0.001162$ \\ 
\hline \\[-1.8ex] 
\end{tabular} 
\end{table} 


\begin{knitrout}
\definecolor{shadecolor}{rgb}{0.969, 0.969, 0.969}\color{fgcolor}\begin{kframe}
\begin{alltt}
\hldef{relative.gain.f2} \hlkwb{<-} \hldef{((var.f2.mc} \hlopt{-} \hldef{var.f2.antithetic)} \hlopt{/} \hldef{var.f2.mc)} \hlopt{*} \hlnum{100}
\end{alltt}
\end{kframe}
\end{knitrout}

\begin{knitrout}
\definecolor{shadecolor}{rgb}{0.969, 0.969, 0.969}\color{fgcolor}\begin{kframe}
\begin{alltt}
\hldef{df10} \hlkwb{<-} \hlkwd{data.frame}\hldef{(}\hlkwc{sample.size} \hldef{= sample.sizes,} \hlkwc{gain} \hldef{= relative.gain.f2)}
\end{alltt}
\end{kframe}
\end{knitrout}


% Table created by stargazer v.5.2.3 by Marek Hlavac, Social Policy Institute. E-mail: marek.hlavac at gmail.com
% Date and time: Tue, Apr 07, 2026 - 22:28:27
\begin{table}[!htbp] \centering 
  \caption{Percentage Variance Reduction in estimation of E[f2(X)]} 
  \label{Table 10} 
\begin{tabular}{@{\extracolsep{5pt}} cc} 
\\[-1.8ex]\hline 
\hline \\[-1.8ex] 
sample.size & gain \\ 
\hline \\[-1.8ex] 
$30$ & $87.182$ \\ 
$80$ & $87.769$ \\ 
$100$ & $86.884$ \\ 
$150$ & $87.137$ \\ 
\hline \\[-1.8ex] 
\end{tabular} 
\end{table} 


\smallpencil \hspace{0.1cm} {\setlength{\spaceskip}{1em plus 0.5em minus 0.5em} \fontsize{17}{20}\myfont We achieve significant reduction in variance of the estimates by using antithetic variable technique. \par}

\end{enumerate}



\end{document}
